% Part: sets-functions-relations
% Chapter: size-of-sets
% Section: zig-zag
% Hindi translation (hi-Deva-IN) of Open Logic commit 9620cc73f9c8e0ad003c514a5d3748f29611c4c0.

\documentclass[../../../include/open-logic-section]{subfiles}

\begin{document}

\olfileid[hi]{sfr}{siz}{zigzag}

\olsection{कैंटर की आड़ी-तिरछी विधि}

\begin{explain}
हम कुछ ``सरल'' गणनों पर पहले ही विचार कर चुके हैं। अब हम कुछ अधिक कठिन विषय
पर विचार करेंगे। प्राकृत संख्याओं के युग्मों के उस समुच्चय पर विचार करें
\oliflabeldef{sfr:set:pai:sec}{, जिसे हमने \olref[set][pai]{sec} में इस
प्रकार परिभाषित किया था:}{, जिसकी परिभाषा है:}
\[
\Nat \times \Nat = \Setabs{\tuple{n,m}}{n,m \in \Nat}
\]
हम इन क्रमित युग्मों को निम्न प्रकार एक \emph{सरणी} में व्यवस्थित कर सकते हैं:
\[
\begin{array}{ c | c | c | c | c | c}
& \mathbf 0 & \mathbf 1 & \mathbf 2 & \mathbf 3 & \dots \\
\hline
\mathbf 0 & \tuple{0,0} & \tuple{0,1} & \tuple{0,2} & \tuple{0,3} & \dots \\
\hline
\mathbf 1 & \tuple{1,0} & \tuple{1,1} & \tuple{1,2} & \tuple{1,3} & \dots \\
\hline
\mathbf 2 & \tuple{2,0} & \tuple{2,1} & \tuple{2,2} & \tuple{2,3} & \dots \\
\hline
\mathbf 3 & \tuple{3,0} & \tuple{3,1} & \tuple{3,2} & \tuple{3,3} & \dots \\
\hline
\vdots & \vdots & \vdots & \vdots & \vdots & \ddots\\
\end{array}
\]
स्पष्ट है कि $\Nat \times \Nat$ का प्रत्येक क्रमित युग्म सरणी में ठीक एक बार
आएगा। विशेष रूप से, $\tuple{n,m}$, $n$वीं पंक्ति और $m$वें स्तंभ में आएगा।
पर ऐसी सरणी के अवयवों को ``एकविमीय'' सूची में कैसे व्यवस्थित करें? नीचे की
सरणी का प्रतिरूप ऐसा करने का एक ढंग दिखाता है (यद्यपि निस्संदेह कई अन्य विकल्प
भी हैं):
\[
\begin{array}{ c | c | c | c | c | c | c}
& \mathbf 0 & \mathbf 1 & \mathbf 2 & \mathbf 3 & \mathbf 4 &\dots \\
\hline
\mathbf 0 & 0  & 1& 3 & 6& 10 &\ldots \\
\hline
\mathbf 1 &2 & 4& 7 & 11 & \dots &\ldots \\
\hline
\mathbf 2 & 5 & 8 & 12 & \ldots & \dots&\ldots \\
\hline
\mathbf 3 & 9 & 13 & \ldots & \ldots & \dots & \ldots \\
\hline
\mathbf 4 & 14 & \ldots & \ldots & \ldots & \dots & \ldots \\
\hline
\vdots & \vdots & \vdots & \vdots & \vdots&\ldots & \ddots\\
\end{array}
\]\noindent
यह प्रतिरूप \emph{कैंटर की आड़ी-तिरछी विधि} कहलाता है। यह
$\Nat \times \Nat$ का गणन इस प्रकार करता है:
\[
\tuple{0,0}, \tuple{0,1}, \tuple{1,0}, \tuple{0,2}, \tuple{1,1},
\tuple{2,0}, \tuple{0,3}, \tuple{1,2}, \tuple{2,1}, \tuple{3,0}, \dots
\]
इससे निम्नलिखित स्थापित होता है:
\end{explain}

\begin{prop}\ollabel{natsquaredenumerable}
$\Nat \times \Nat$ !!{enumerable} है।
\end{prop}

\begin{proof}
फलन $f \colon \Nat \to \Nat\times\Nat$ प्रत्येक $k \in \Nat$ को उस
युग्म $\tuple{n,m} \in \Nat \times \Nat$ पर भेजे जिसके लिए $k$, कैंटर की
आड़ी-तिरछी सरणी की $n$वीं पंक्ति और $m$वें स्तंभ में लिखा मान है। 
\end{proof}

\begin{explain}
यह तकनीक बड़े सुगम ढंग से व्यापक भी होती है। उदाहरणार्थ, हम इसका उपयोग प्राकृत
संख्याओं के क्रमित त्रिकों के समुच्चय का गणन करने में कर सकते हैं, अर्थात:
\[
\Nat \times \Nat \times \Nat = \Setabs{\tuple{n,m,k}}{n,m,k \in \Nat}
\]
हम $\Nat \times \Nat \times \Nat$ को $\Nat \times \Nat$ और $\Nat$ का
कार्तीय गुणनफल मानते हैं, अर्थात,
\[
\Nat^3 = (\Nat \times \Nat) \times \Nat =
\Setabs{\tuple{\tuple{n,m},k}}{n, m, k
  \in \Nat }
\]
और इस प्रकार हम $\Nat^3$ का गणन ऐसी सरणी से कर सकते हैं जिसके एक अक्ष पर
$\Nat$ का गणन और दूसरे अक्ष पर $\Nat^2$ का गणन अंकित हो:
\[
\begin{array}{ c | c | c | c | c | c}
& \mathbf 0 & \mathbf 1 & \mathbf 2 & \mathbf 3 & \dots \\
\hline
\mathbf{\tuple{0,0}} & \tuple{0,0,0} & \tuple{0,0,1} & \tuple{0,0,2} & \tuple{0,0,3} & \dots \\
\hline
\mathbf{\tuple{0,1}} & \tuple{0,1,0} & \tuple{0,1,1} & \tuple{0,1,2} & \tuple{0,1,3} & \dots \\
\hline
\mathbf{\tuple{1,0}} & \tuple{1,0,0} & \tuple{1,0,1} & \tuple{1,0,2} & \tuple{1,0,3} & \dots \\
\hline
\mathbf{\tuple{0,2}} & \tuple{0,2,0} & \tuple{0,2,1} & \tuple{0,2,2} & \tuple{0,2,3} & \dots\\
\hline
\vdots & \vdots & \vdots & \vdots & \vdots & \ddots \\
\end{array}
\]
इस प्रकार, कैंटर की आड़ी-तिरछी विधि जैसी विधि से हम इसी तरह~$\Nat^3$ का गणन
प्राप्त कर सकते हैं। और हम आगे बढ़ते हुए $\Nat^n$ के गणन प्राप्त कर सकते हैं,
जहाँ $n$ कोई भी प्राकृत संख्या हो। अतः:
\end{explain}

\begin{prop}
$\Nat^n$ प्रत्येक $n \in \Nat$ के लिए !!{enumerable} है।
\end{prop}

\begin{prob}\label{sfr:siz:zigzag:prob:posint-n}
दिखाइए कि $(\PosInt)^n$ प्रत्येक $n \in \Nat$ के लिए !!{enumerable} है।
\end{prob}

\begin{prob}\label{sfr:siz:zigzag:prob:posint-star}
  दिखाइए कि $(\PosInt)^*$ !!{enumerable} है। आप \cref{sfr:siz:zigzag:prob:posint-n}
  को मान सकते हैं।
\end{prob}

\end{document}
